
\textbf{Proposition 5.2.} \uline{On a sur} @@M0@@ \uline{la suite exacte} :
@@M1@@
\uline{i}.\uline{e}. \uline{on a un isomorphisme}
@@M2@@

On considère le fibré vectoriel sur @@M3@@ défini par @@M4@@, et on note @@M5@@ le complémentaire de la section nulle de ce fibré. On a le diagramme commutatif :
@@M6@@
Le morphisme @@M7@@ identifie @@M8@@ au @@M9@@-schéma affine défini par le faisceau d'algèbres @@M10@@ ; @@M11@@ est une @@M12@@-Algèbre graduée.

Comme @@M13@@ sont des morphismes lisses, on peut écrire la suite exacte des faisceaux de différentielles relatives ; comme @@M14@@ est affine sur @@M15@@, on peut écrire cette suite en termes de @@M16@@-Modules; on obtient ainsi :
@@M17@@

Or cette suite exacte est une suite exacte graduée de @@M18@@-Modules gradués. En effet, elle résulte du diagramme commutatif :
@@M19@@
où les flèches sont des homomorphismes d'algèbres, compatibles avec la graduation de @@M20@@ (graduation définie par le degré par rapport à l'indéterminée @@M21@@), les deux autres algèbres étant supposées réduites
